% id: 140-19
% book: 베이직쎈 미적분1 · 09 정적분의 활용
% section: 개념 57 두 곡선 사이의 넓이
% passage: 다음 곡선과 직선으로 둘러싸인 도형의 넓이를 구하시오.
% figure: tikz:fig-140-19
% answer: 
% answer_source: 
% variant_level: 0 (원본 전사)
% 이 파일은 build.mjs 가 items.json 에서 만든다. 고칠 때는 items.json 을 고친다.
\smash{\raisebox{1.5mm}{\dmkichul{베이직쎈 미적분1 140쪽 19번}}}
\noindent\begin{minipage}[t]{0.58\linewidth}\vspace{0pt}\raggedright $y=x^2$, $y=2x+3$\end{minipage}\hfill\begin{minipage}[t]{0.4\linewidth}\vspace{0pt}\centering \resizebox{\ifdim\width>\linewidth\linewidth\else\width\fi}{!}{% 140-19(140쪽) — 곡선 $y=x^2$ 과 직선 $y=2x+3$ 으로 둘러싸인 색칠한 부분($x=-1$ 에서 $x=3$ 까지). 스캔 화질이 낮아 TikZ 로 다시 그림.
% 원문에 있는 것만 옮긴다: 포물선 · 직선 · 색칠 · $x=3$ 의 안내 점선 · 지시선 · 라벨 y=x^2, y=2x+3, -1, O, 3, x, y. 원문에 없는 눈금·넓이 값은 적지 않는다.
\begin{tikzpicture}[x=0.19cm, y=0.16cm, line width=0.6pt]
  \fill[gray!25] plot[domain=-1:3, samples=80, smooth] (\x, {2*(\x)+3}) -- plot[domain=3:-1, samples=80, smooth] (\x, {(\x)*(\x)}) -- cycle;
  \draw[->, >=stealth, line width=0.7pt] (-3.8,0) -- (4.7,0) node[below=1pt, font=\small] {$x$};
  \draw[->, >=stealth, line width=0.7pt] (0,-1.7) -- (0,13.3) node[left=1pt, font=\small] {$y$};
  \draw[densely dotted, line width=0.4pt] (3,0) -- (3,9);
  \draw[domain=-3.35:3.3, samples=120, smooth] plot (\x, {(\x)*(\x)});
  \draw (-2.25,-1.5) -- (3.6,10.2);
  \draw[->, >=stealth, line width=0.4pt] (1.8,11.7) .. controls (2.5,11.1) and (2.8,10.2) .. (2.82,8.1);
  \node[anchor=west, inner sep=1pt, font=\small] at (0.55,12.6) {$y=x^2$};
  \node[anchor=west, inner sep=1pt, font=\small] at (3.5,9.9) {$y=2x+3$};
  \node[below=1pt, font=\small] at (-1.2,0) {$-1$};
  \node[below right=-1pt and -2pt, font=\small] at (0,0) {$\pt{O}$};
  \node[below=1pt, font=\small] at (3,0) {$3$};
\end{tikzpicture}
}\end{minipage}\par
\dmhint{곡선 $y=x^2$과 직선 $y=2x+3$의 교점의 $x$좌표는 $x^2=2x+3$에서\\ \quad $x^2-2x-3=0$\\ \quad $(x+\blank[1.1em])(x-3)=0$\\ \quad $\therefore x=\blank[1.1em]$ 또는 $x=3$\\ 따라서 구하는 넓이는\\ \quad $\displaystyle\int_{\text{\scriptsize\setlength{\fboxsep}{1.5pt}\framebox[1.5em]{\rule{0pt}{1.6ex}}}}^{3} (2x+3-\blank[1.1em])\mkern3mu dx$\\ \quad $\displaystyle =\Big[x^2+\blank[1.1em]-\dfrac{1}{3}x^3\Big]_{\text{\scriptsize\setlength{\fboxsep}{1.5pt}\framebox[1.5em]{\rule{0pt}{1.6ex}}}}^{3}=\blank$}
